\chapter{Background}\label{background}

This chapter presents the technical foundations underlying this project.
Section~\ref{sec:gnn} introduces Graph Neural Networks (GNNs) and the
message-passing paradigm. Section~\ref{sec:gnn-ehr} reviews the application of
GNNs to Electronic Health Records (EHRs) for clinical risk prediction.
Section~\ref{sec:what-is-xai} introduces Explainable AI and the
interpretability--performance trade-off in the clinical context.
Section~\ref{sec:xai-gnn} discusses explainability in GNNs and highlights the
limitations of post-hoc approaches. Finally, Section~\ref{sec:proto} introduces
prototype learning as an inherently interpretable alternative.

%===============================================================
\section{Graph Neural Networks}
\label{sec:gnn}
%===============================================================

A graph is defined as $G = (V, E)$, where $V$ denotes the set of nodes and
$E \subseteq V \times V$ represents the set of edges. Each node $v \in V$ is
associated with a feature vector $\mathbf{x}_v \in \mathbb{R}^d$, and edges may
also carry features. Graph Neural Networks (GNNs) are deep learning models
designed to learn representations of nodes, edges, or entire graphs by
exploiting the underlying graph structure~\cite{yuan2022explainability}.

The core mechanism of GNNs is \emph{message passing}. At layer $k$, each node
updates its representation by aggregating information from its neighbours
$\mathcal{N}(v)$:
%
\begin{equation}
\mathbf{h}_v^{(k+1)} = \sigma\!\left(
  \sum_{u \in \mathcal{N}(v) \cup \{v\}}
  \tilde{\mathbf{A}}_{uv}\, \mathbf{W}^{(k)} \mathbf{h}_u^{(k)}
\right),
\end{equation}
%
where $\mathbf{W}^{(k)}$ is a learnable weight matrix, $\tilde{\mathbf{A}}$ is
the normalised adjacency matrix including self-loops, and $\sigma(\cdot)$ is a
non-linear activation function~\cite{kipf2017gcn}. After $L$ layers, node
embeddings capture information from their $L$-hop neighbourhoods. For
graph-level tasks, a readout function (e.g.\ mean, max, or sum pooling)
aggregates node embeddings into a single graph representation $\mathbf{h}_G$.

In this work we build on two widely used GNN architectures. The Graph
Convolutional Network (GCN)~\cite{kipf2017gcn} applies symmetric normalisation
to the adjacency matrix during aggregation and serves as the encoder of our
model. The Graph Attention Network (GAT)~\cite{velickovic2018gat} instead
learns attention weights that assign different importance to neighbouring
nodes, which is useful in clinical settings where the relevance of relations
between entities may vary~\cite{ossboll2024gnn}; it is implemented as an
alternative encoder for comparison.

%===============================================================
\section{Graph Neural Networks for Electronic Health Records}
\label{sec:gnn-ehr}
%===============================================================

EHRs are inherently relational: a single emergency-department visit links a
patient to diagnoses, medications, laboratory results, procedures, and vital
signs, and a patient accumulates many such visits over time. This structure
makes GNNs a natural fit for clinical prediction, and a growing body of work
applies them to tasks such as mortality prediction, readmission, and
diagnosis~\cite{ossboll2024gnn}.

Two graph-construction paradigms dominate. In the first, \emph{population} (or
\emph{patient-similarity}) graphs, each patient is a node and edges connect
clinically similar patients, so that label information can propagate between
comparable cases. In the second, \emph{intra-patient} graphs, the medical
entities of a visit (diagnoses, medications, vitals) become nodes and edges
encode their clinical relationships; prediction is then a graph-classification
task with one graph per patient. The latter yields a topology that is
clinically meaningful rather than a pure similarity heuristic, and it is the
formulation adopted in this work (Chapter~\ref{architecture}).

Several methods build on these paradigms. Temporal approaches such as DynaGraph
model the sequence of visits as a dynamic graph to capture disease
progression~\cite{mesinovic2026dynagraph}. GraphCare follows a knowledge-graph
route and serves as the main GNN comparison baseline in this project; it is
described in Section~\ref{sec:graphcare}. Most of these models remain black
boxes: they predict accurately but do not expose which entities or relations
drove a decision. ProtoEHR is the closest prior work to ours: it brings
prototype-based, case-based reasoning to hierarchical EHR data, but it does not
identify which parts of the patient graph drive each prediction, nor does it
compare its explanations against post-hoc methods~\cite{cai2025protoehr}.

%===============================================================
\section{Explainable Artificial Intelligence (XAI)}
\label{sec:what-is-xai}
%===============================================================

As Artificial Intelligence (AI) systems are increasingly integrated into
critical domains like healthcare, the ``black-box'' nature of deep learning
models has become a primary barrier to adoption. Explainable AI (XAI) refers to
machine learning techniques designed to make the predictions of AI models
understandable to human experts~\cite{yuan2022explainability}.

\subsection{The Interpretability--Performance Trade-off}\label{sec:tradeoff}

Historically, it was assumed that there exists a trade-off between model
complexity (performance) and interpretability. Simple models, such as linear
regression or decision trees, are more intuitive and interpretable but often
fail to capture the complex, non-linear dependencies found in multi-modal EHR
data. Conversely, GNNs hold high predictive power but fail to provide insight
into \emph{how and why} a specific outcome was reached. XAI seeks to resolve
this either by creating post-hoc tools to probe complex models or by designing
``interpretable-by-design'' architectures.

\subsection{XAI in the Clinical Context}

In healthcare, interpretability has become a requirement for safety and
accountability. Clinical XAI must satisfy the ``right to explanation'' under
legal frameworks like the GDPR and provide practitioners with actionable and
reliable evidence. For instance, an explainer that highlights a patient's age
as the primary risk factor for a disease is less clinically useful than one
that highlights a specific combination of rising creatinine levels and falling
blood pressure.

%===============================================================
\section{Explainability in Graph Neural Networks}
\label{sec:xai-gnn}
%===============================================================

Despite their strong predictive performance, GNNs are often considered
black-box models. Existing explainability methods can be broadly categorised
into instance-level and model-level approaches~\cite{yuan2022explainability}.
Instance-level methods aim to explain individual predictions and can be further
divided into gradient-based, perturbation-based, decomposition-based, and
surrogate-based techniques.

A common characteristic of these methods is that they are \emph{post-hoc}: they
attempt to explain a trained model after the prediction has been made.
Typically, they identify a subgraph that maximises the mutual information with
the model's output. However, this design introduces two key limitations. First,
explanations are not guaranteed to reflect the true reasoning of the model,
since the explanation mechanism is separate from the prediction
process~\cite{zhang2022protgnn}. Second, large-scale evaluations such as
GraphXAI show that many post-hoc methods struggle with large or complex
explanations and may fail to preserve fairness properties, both of which are
critical in clinical applications~\cite{agarwal2023graphxai}.

To evaluate explainability methods, the literature relies on metrics such as
fidelity, accuracy, stability, and sparsity~\cite{agarwal2023graphxai}. Because
the clinical data in this project has no ground-truth importance masks,
explanation quality is assessed primarily through sparsity, inter-explainer
agreement, and, optionally, fidelity, while ground-truth-dependent metrics
such as explanation accuracy are not applicable
(Chapter~\ref{implementation}).

%===============================================================
\section{Prototype Learning and Case-Based Reasoning}
\label{sec:proto}
%===============================================================

Prototype learning is inspired by \emph{case-based reasoning}, where new
instances are classified by comparing them to representative examples. In deep
learning, this idea was popularised by ProtoPNet, which introduces a prototype
layer that stores representative patterns and classifies inputs based on their
similarity to these prototypes~\cite{chen2019protopnet}.

This concept has been extended to graph data through ProtGNN. The model consists
of three components: a GNN encoder $f$ that produces a graph embedding
$\mathbf{h}$, a prototype layer $g_P$ containing $K$ learnable prototype vectors
$\{\mathbf{p}_j\}$, and a final classification layer~\cite{zhang2022protgnn}.
The similarity between an embedding and a prototype is computed as
%
\begin{equation}
\mathrm{sim}(\mathbf{p}_k, \mathbf{h}) = \log\!\left(
  \frac{\lVert \mathbf{p}_k - \mathbf{h}\rVert_2^2 + 1}
       {\lVert \mathbf{p}_k - \mathbf{h}\rVert_2^2 + \epsilon}
\right),
\end{equation}
%
which yields a positive, distance-based similarity score that is maximal when
the embedding coincides with the prototype and decays towards zero as the
distance grows.

The training objective combines a classification loss with regularisation terms
that shape the prototype space: a cluster loss (encouraging embeddings to be
close to same-class prototypes), a separation loss (pushing embeddings away from
other-class prototypes), and, in the original ProtGNN, a diversity loss
(preventing prototype collapse)~\cite{zhang2022protgnn}. The clinical adaptation
used in this project retains the cluster and separation terms but replaces the
diversity term with an $\ell_1$ sparsity penalty on the classification layer
(Chapter~\ref{implementation}).

A key feature of prototype learning is \emph{prototype projection}, where each
prototype is periodically aligned with the closest real training example. In
ProtGNN, this is achieved using Monte Carlo Tree Search (MCTS) to efficiently
explore the space of possible subgraphs~\cite{zhang2022protgnn}. Unlike post-hoc
methods, prototype-based models provide explanations as part of the prediction
process itself, making them inherently interpretable. In the clinical domain,
ProtoEHR extends this idea to hierarchical EHR data, making it the closest prior
work to this project~\cite{cai2025protoehr}.

%===============================================================
\section{Knowledge-Graph GNNs and GraphCare}
\label{sec:graphcare}
%===============================================================

Prototype models represent one route to interpretable EHR prediction. A second
route enriches the patient graph with external medical knowledge before a GNN is
applied. GraphCare is the representative method of this route and serves as the
main GNN comparison baseline in this project~\cite{jiang2024graphcare}.

GraphCare builds a personalised knowledge graph for each patient. Rather than
relying only on the relations present in the dataset, it extracts relational
knowledge about medical concepts from large language models and from external
biomedical knowledge bases, and attaches it to the patient's diagnoses,
procedures, and medications. The result is a patient-specific graph that carries
context a raw record does not contain.

Prediction is performed by a Bi-attention Augmented (BAT) graph neural network.
The bi-attention mechanism weights both the nodes of the personalised graph and
the visits over time, so the model can focus on the concepts and moments most
relevant to a given prediction. In the original work, GraphCare was evaluated on
MIMIC-III and MIMIC-IV across mortality prediction, readmission, length of stay,
and drug recommendation, and reported its largest gains in settings with limited
training data.

Two properties matter for this project. GraphCare and the proposed prototype
model share the same graph-classification setting, which makes a controlled
comparison on the same split meaningful. Their interpretability differs in kind:
GraphCare exposes attention weights over the personalised graph, whereas the
prototype model explains a prediction through similarity to learned cases and
through node attribution over the patient graph. This difference is the reason
GraphCare is used as the reference point in the evaluation
(Chapter~\ref{evaluation}).

\section{Method 3}
\label{sec:method3}